\documentclass[11pt,a4paper]{article}
\usepackage[T1]{fontenc}
\usepackage{isabelle,isabellesym}
\usepackage{amsmath}
\usepackage{amssymb}
\usepackage{url}

\isabellestyle{it}

\begin{document}

\title{Structural Metatheory for Finitary One-Dimensional Higher Inductive Types}
\author{Arthur Freitas Ramos \and David Barros Hulak \and
  Ruy J. G. B. de Queiroz}
\maketitle

\begin{abstract}
This stand-alone entry develops the structural metatheory of a finitary,
one-dimensional, nondependent calculus of higher inductive types directly in
Isabelle/HOL. It has no dependency on a homotopy type theory object logic. Its
terms, raw typing, and reduction are inductive, and its signature environment
uniformly describes ordinary inductive types and higher inductive types.

A signature uniformly declares finite lists of point and path constructors.
Each endpoint expression is checked against constructor arities and argument
types, and instantiating a checked endpoint is proved to preserve typing.
Free dimension support is defined structurally and proved to be bounded by the
ambient dimension context; substitution at a fresh dimension is therefore the
identity on every well-typed term.
Path constructors are \emph{cubical-style}: a constructor takes a dimension
argument, its faces are obtained by substituting the two constant dimensions,
and its eliminator computes judgmentally. A path method is accepted only when
both typed faces are joinable with the eliminations of the corresponding
declared endpoints, and hence convertible to them. Structural elimination of 
a checked endpoint is proved
well typed. The resulting boundary-computation square states that eliminating
either declared face is convertible to taking the same face of the path method
after computation. This prevents the interval counterexample in which a
constant segment method disagrees with one of the point methods.

Boundary admissibility is no longer left only as a quantification over every
context and every typed constructor instance. A finite open-term criterion
checks the two faces of each path method once, under canonical variables
appended for that constructor's arguments. A soundness theorem for simultaneous
substitution transports these finitely many open obligations to all typed
instances. Strong normalization and confluence are also packaged as a canonical
normal-form function, and typed joinability is proved equivalent to equality of
canonical normal forms. Consequently, the finite boundary criterion has an
equivalent normal-form formulation. The canonical function uses Hilbert choice,
so this is a theorem-level comparison criterion rather than an executable
normalizer. Separately, the hereditary checked judgment remains explicitly a
semantic safety invariant: every reduct must satisfy the local structural
predicate.

The structural metatheory proves dimension-context weakening, de Bruijn
weakening, capture-avoiding term substitution, dimension substitution, and
one-step and multi-step preservation of raw typing and hereditary safety.
Every raw well-typed term is strongly normalizing. Local confluence on regular
terms, together with typing regularity and strong normalization, yields
Church--Rosser confluence and unique normal forms. Closed checked terms
therefore reach canonical normal forms; at higher inductive type these are
point constructors with canonical arguments. A signature with no point
constructors has no closed checked inhabitant.
\end{abstract}

\tableofcontents

\section{Scope}

The calculus has two type formers, function spaces and higher inductive types
drawn from the signature environment. Constructor arguments are either
parameters of fixed higher inductive type or recursive occurrences. The schema
covers the circle, interval, suspensions, and the indiscrete paths of
propositional truncation. It does not express general quotients by a relation
family whose witness type depends on earlier arguments.

Infinitary constructors, dependent constructor arguments, indexed families,
dependent eliminators, higher path constructors, dimension abstraction,
composition, and Kan operations are outside the scope. Thus ``cubical-style''
describes the dimension-indexed path syntax and judgmental path computation;
this is not a formalization of cubical type theory. The metatheory does prove
strong normalization and Church--Rosser confluence for its stated finitary,
one-dimensional, nondependent calculus; it does not claim those results for
the richer systems excluded above.

\section{Related work and position}

{\sloppy
The general shape of the schema follows the informal descriptions of higher
inductive definitions in the Homotopy Type Theory book \cite{HoTTBook2013} and
the schemata studied by Lumsdaine and Shulman \cite{LumsdaineShulman2020} and
by Sojakova \cite{Sojakova2015}. Kaposi and Kov\'acs give a substantially more
expressive theory of signatures for higher inductive-inductive types, formalize
its syntactic translations in Agda, and generate induction principles
\cite{KaposiKovacs2020}. The present calculus is narrower---nondependent,
one-sorted at each signature index, and one-dimensional---but proves an
operational metatheory not supplied by that signature translation: explicit
reduction, subject reduction, strong normalization, Church--Rosser, unique
normal forms, and consistency.

The closest Isabelle-specific precedent is Chen's Isabelle/HoTT
\cite{Chen2021}, the first implementation of homotopy type theory in Isabelle.
It uses Isabelle's logical-framework infrastructure for intensional dependent
type theory, type checking, elaboration, Isar proofs, and foundational HoTT
results. Individual higher inductive constants and rules can be postulated in
such an object logic, but Isabelle/HoTT does not supply a uniform HIT signature
schema, an open boundary-admissibility criterion, or the operational
metatheory proved here. This entry instead reifies its calculus as datatypes and
relations in a stand-alone Isabelle/HOL session.

Veltri and van der Weide define internal signatures for finitary 1-truncated
higher inductive types and construct their algebras from groupoid quotients in
UniMath \cite{VeltriVanDerWeide2021}. Their result is an internal semantic
construction with initial-algebra content. This entry instead studies an
external first-order term calculus and verifies its operational properties in
Isabelle; it neither constructs semantic models nor claims the expressive
generality of their signatures.

The judgmental treatment of path constructors is inspired by cubical type
theory \cite{CohenCoquandHuberMoertberg2018}. Cavallo and Harper provide a
general indexed cubical inductive schema with computational higher inductive
types and canonicity \cite{CavalloHarper2019}; Huber proves canonicity for
cubical type theory \cite{Huber2019}; and Sterling and Angiuli prove
normalization and decidability of judgmental equality for Cartesian cubical
type theory \cite{SterlingAngiuli2021}. In contrast, the present system has no
composition, Kan operations, dimension abstraction, indexed families, or
higher path constructors. Its distinct contribution is a compact
one-dimensional operational calculus with explicit boundary-coherence and
dimension-capture analyses, together with machine-checked strong normalization
and Church--Rosser proofs.
\par}

Dimension variables are natural numbers, and the index~$0$ is reserved for the
dimension bound by the path methods of an eliminator. Ambient dimension
contexts therefore consist of strictly positive indices, and dimension
substitution never has to shift.

\section{The formal development}

\input{session}

\section{AI assistance}

AI assistance was used for proof engineering. The final definitions,
statements, and proofs are checked by Isabelle.

\bibliographystyle{abbrv}
\bibliography{root}

\end{document}
